% ============================================================
\chapter{Direct Methods for Linear Systems}
\label{ch:direct}
% ============================================================

\section{Motivating example}
An electrical circuit with $n$ nodes produces a system $Ax=b$ of $n$ equations. For $n=1000$, a systematic and efficient method is needed.

\section{Gaussian elimination and LU factorisation}

\begin{definition}[LU factorisation]\index{LU factorisation}\index{LU}
Find $L$ lower triangular (with $\ell_{ii}=1$) and $U$ upper triangular such that $A=LU$. Then $Ax=b$ is solved by $Ly=b$ (forward substitution) followed by $Ux=y$ (back substitution).
\end{definition}

\begin{algorithm}[H]
\caption{LU factorisation (without pivoting)}\label{alg:lu}
\KwInput{Matrix $A\in\R^{n\times n}$}
\KwOutput{$L$, $U$ such that $A=LU$}
\For{$k=1,\ldots,n-1$}{
  \For{$i=k+1,\ldots,n$}{
    $\ell_{ik} \gets a_{ik}/a_{kk}$\;
    \For{$j=k+1,\ldots,n$}{
      $a_{ij} \gets a_{ij}-\ell_{ik}\cdot a_{kj}$\;
    }
  }
}
$U \gets$ upper triangular part of $A$\;
\end{algorithm}

\begin{theorem}[Existence of LU]
If all leading principal submatrices of $A$ are invertible, then the LU factorisation exists and is unique.
\end{theorem}

\begin{theorem}[Cost]
The LU factorisation costs $\frac{2n^3}{3}+\BigO(n^2)$ floating-point operations.
\end{theorem}

\begin{proof}
Step $k$ performs $(n-k)^2$ multiplications and subtractions. Total: $\sum_{k=1}^{n-1}2(n-k)^2=2\sum_{j=1}^{n-1}j^2=\frac{2n^3}{3}+\BigO(n^2)$.
\end{proof}

\section{Partial pivoting}\index{partial pivoting}

\begin{warning}
Without pivoting, the algorithm can be unstable (division by a small pivot). \textbf{Partial pivoting} permutes rows to maximise $|a_{kk}|$.
\end{warning}

One obtains $PA=LU$ where $P$ is a permutation matrix.

\begin{theorem}[Stability]
With partial pivoting, Gaussian elimination is \textbf{backward stable}:
\[
(A+\Delta A)\hat{x}=b, \quad \|\Delta A\|\leq g(n)\,\eps_{\mathrm{mach}}\,\|A\|,
\]
where $g(n)$ is the growth factor (bounded by $2^{n-1}$ in theory, rarely exceeded in practice).
\end{theorem}

\section{Cholesky factorisation}\index{Cholesky}

\begin{theorem}[Cholesky factorisation]
If $A$ is \textbf{symmetric positive definite} (SPD), then there exists a unique lower triangular matrix $L$ with positive diagonal entries such that $A=LL^\top$.
\end{theorem}

\begin{proof}
By induction on $n$. Write $A=\begin{pmatrix}a_{11}&\mathbf{v}^\top\\\mathbf{v}&A'\end{pmatrix}$. Since $A$ is SPD, $a_{11}>0$. Set $\ell_{11}=\sqrt{a_{11}}$, $\mathbf{l}=\mathbf{v}/\ell_{11}$. Then $A'- \mathbf{l}\mathbf{l}^\top$ is SPD of size $(n-1)\times(n-1)$ and we apply the induction hypothesis.
\end{proof}

\begin{algorithm}[H]
\caption{Cholesky factorisation}\label{alg:cholesky}
\KwInput{$A$ SPD, $n\times n$}
\KwOutput{$L$ such that $A=LL^\top$}
\For{$j=1,\ldots,n$}{
  $\ell_{jj}\gets\sqrt{a_{jj}-\sum_{k=1}^{j-1}\ell_{jk}^2}$\;
  \For{$i=j+1,\ldots,n$}{
    $\ell_{ij}\gets\frac{1}{\ell_{jj}}\left(a_{ij}-\sum_{k=1}^{j-1}\ell_{ik}\ell_{jk}\right)$\;
  }
}
\end{algorithm}

\begin{theorem}[Cost of Cholesky]
$\frac{n^3}{3}+\BigO(n^2)$: half the cost of LU.
\end{theorem}

\section{Numerical example}

$A=\begin{pmatrix}4&2\\2&5\end{pmatrix}$, $b=\begin{pmatrix}8\\9\end{pmatrix}$. Cholesky: $L=\begin{pmatrix}2&0\\1&2\end{pmatrix}$.

Forward substitution $Ly=b$: $y_1=4$, $y_2=(9-1\cdot4)/2=2.5$. Back substitution $L^\top x=y$: $x_2=2.5/2=1.25$, $x_1=(4-1\cdot 1.25)/2=1.375$.

\section{Implementation}

\begin{codeblock}[title=\textbf{Python}]
\begin{minted}{python}
import numpy as np
from scipy.linalg import lu, cholesky, solve_triangular

# LU
A = np.array([[2., 1, 1], [4, 3, 3], [8, 7, 9]])
b = np.array([4., 10, 24])
P, L, U = lu(A)
y = solve_triangular(L, P @ b, lower=True)
x = solve_triangular(U, y)
print(f"LU: x = {x}")

# Cholesky
A_spd = np.array([[4., 2], [2, 5]])
b_spd = np.array([8., 9])
L_chol = cholesky(A_spd, lower=True)
y = solve_triangular(L_chol, b_spd, lower=True)
x = solve_triangular(L_chol.T, y)
print(f"Cholesky: x = {x}")
print(f"Cout LU (n=100): ~{2*100**3//3:.0f} flops")
\end{minted}
\end{codeblock}

\begin{codeblock}[title=\textbf{Julia}]
\begin{minted}{julia}
using LinearAlgebra
A = [2.0 1 1; 4 3 3; 8 7 9]
b = [4.0, 10, 24]
F = lu(A)
x = F \ b
println("LU: x = $x")

A_spd = [4.0 2; 2 5]
b_spd = [8.0, 9]
C = cholesky(A_spd)
x = C \ b_spd
println("Cholesky: x = $x")
\end{minted}
\end{codeblock}

\section{Exercises}

\begin{exercise}[$\star$]
Factorise $A=\begin{pmatrix}1&2&3\\2&8&14\\3&14&34\end{pmatrix}$ into $LU$ by hand. Verify.
\end{exercise}

\begin{exercise}[$\star\star$]
Show that if $A$ is SPD, all pivots in Gaussian elimination are positive (no pivoting needed).
\end{exercise}

\begin{exercise}[$\star\star$]
Compare the computation times of LU vs.\ Cholesky for $n=100,500,1000,2000$ on random SPD matrices. Verify the ratio $\sim 2$.
\end{exercise}

\begin{exercise}[$\star\star\star$ -- Project]
Implement the banded LU factorisation for tridiagonal matrices (Thomas algorithm). Apply to the discretisation of $-u''=f$.
\end{exercise}

\begin{keyformulas}
\begin{align*}
A &= LU \quad (\text{cost } \tfrac{2n^3}{3}) \\
A &= LL^\top \quad (\text{Cholesky, cost } \tfrac{n^3}{3}) \\
PA &= LU \quad (\text{partial pivoting}) \\
\cond(A) &= \|A\|\cdot\|A^{-1}\|, \quad \frac{\|\delta x\|}{\|x\|}\leq\cond(A)\frac{\|\delta b\|}{\|b\|}
\end{align*}
\end{keyformulas}
